\subsection{A significant revision: retaining uncertain evidence}
A controlled comparison shows why evidence bounds belong in the decision.
On the same 160 utility entries, compare the midpoint completion $z_{ij}=(\ell_{ij}+h_{ij})/2$ with the retained-interval formulation in Equation~\ref{eq:regret}. Both use identical weights, anchors and TOPSIS scores. The midpoint is a declared diagnostic scenario, rather than an estimate of missing observations or a reconstruction of the team's historical first model.

The revision retains 48 non-point entries and selects the city with the smallest worst-case shortfall. Zurich remains the recommendation, but 13 cities change decision rank under fixed weights; Kendall's $\tau_b=0.583$. With equal weights, 15 ranks change and $\tau_b=0.650$.

For each target city $c$, construct a separate interior completion: set every uncertain utility to $\ell_{ij}+\alpha_i(h_{ij}-\ell_{ij})$, with $\alpha_c=0.95$ and $\alpha_i=0.05$ for $i\ne c$. All 112 point values remain fixed and all 48 uncertain values lie strictly inside their intervals. Each target ranks first under both baseline and equal weights. Across the sixteen targets, the smallest lead is 4.974 score points under baseline weights and 5.323 under equal weights. These separate completions establish feasibility in the declared utility rectangle; they carry no estimated probabilities and do not establish joint physical attainability.

Retaining intervals changes the decision order while preserving sixteen possible champions. The \TotalWidthReduction\% reduction in interval width comes from admitted data, separately from this formulation comparison.
